\documentclass[11pt,a4paper]{article}
\usepackage[utf8]{inputenc}
\usepackage[T1]{fontenc}
\IfFileExists{polish.ldf}{\usepackage[polish]{babel}}{}
\IfFileExists{lmodern.sty}{\usepackage{lmodern}}{}
\usepackage[margin=26mm]{geometry}
\usepackage{amsmath}
\usepackage[authoryear,round]{natbib}
\usepackage[hidelinks]{hyperref}
\title{Trafność punktowego pomiaru spojrzenia z kamery Android: kalibracja i powtarzalność}
\author{Mateusz Pomianek}
\date{MSC-EYE-P01 | Protokół prospektywny 1.0.0 | 27 listopada 2026}
\begin{document}
\maketitle
\begin{abstract}
Protokół ocenia punktowy pomiar spojrzenia na ekranie Android po kalibracji indywidualnej. Rozdziela błąd względem celu fiksacji, dyspersję wewnątrz fiksacji, dostępność oszacowań i powtarzalność między sesjami. Pierwszorzędowy estymand dotyczy błędu na rozłącznym zbiorze celów w kontrolowanym warunku. Obliczenia geometrii zweryfikowano na przypadkach syntetycznych. Badanie z udziałem osób wymaga działającego toru pomiarowego i dopuszczenia protokołu.
\end{abstract}
\section{Wprowadzenie}
Estymacja spojrzenia z obrazu kamery mobilnej zależy między innymi od kalibracji i warunków rejestracji \citep{krafka}. Celem jest ilościowa ocena konkretnej, przypiętej implementacji Android w zadaniu fiksacji celów na ekranie. P01 bada lokalizację ciągłą. P02 dotyczy wyboru wiarygodnych oszacowań po zmianie warunków, a P03 dotyczy klasyfikacji regionów w grze. Położenie bodźca stanowi cel zadania, a zgodność oszacowania z rzeczywistą osią wzroku wymaga odrębnego porównania z aparaturą referencyjną.

\section{Materiał i metody}
W pierwszej fazie kwalifikuje się rejestrację obrazu, układ współrzędnych, wersję estymatora, kalibrację i eksport. W fazie E01 każda osoba wykonuje dwie wizyty, z ponowną kalibracją podczas każdej z nich. Analiza główna wykorzystuje pierwszy zaplanowany blok po kalibracji z każdej wizyty. Kolejność celów oceny losuje się niezależnie od kolejności kalibracji. Cele oceny mają współrzędne różne od celów dopasowania i nie są wykorzystywane do strojenia. Próby po nieudanej kalibracji otrzymują jawny status. Odległość, ustawienie telefonu, oświetlenie i nominalny zakres pozy głowy są rejestrowane. Badanie odporności rozdziela zmiany odległości, pozy i oświetlenia od warunku bazowego. Wnioski o różnicach sprzętu wymagają osobnego projektu obejmującego powtarzanie warunków na kilku urządzeniach.

Dla celu i estymaty wyznacza się wektory $\mathbf u$ i $\mathbf v$ od przyjętego punktu oka do płaszczyzny ekranu. Błąd kątowy oblicza się jako
\[
 \theta=\frac{180}{\pi}\operatorname{atan2}
 \bigl(\lVert\mathbf u\times\mathbf v\rVert,\;\mathbf u\cdot\mathbf v\bigr).
\]
Geometria wymaga skali poziomej i pionowej, odległości i rzutu punktu oka na ekran. Gdy geometria nie jest znana, raportuje się błąd pikselowy ze wskazaniem układu odniesienia. Opóźnienie i dyspersja są oddzielnymi charakterystykami.

\section{Plan analizy}
Pierwszorzędowy wskaźnik: średnia międzyosobnicza z median kątowego błędu prób względem celu, w warunku bazowym pierwszej wizyty i przy ważnej kalibracji. Błąd pojedynczej próby pochodzi z mediany współrzędnych w oknie fiksacji. Osoby mają równe wagi. Obok wyniku raportuje się odsetek prób z użytecznym oszacowaniem wśród wszystkich zaplanowanych prób, także przy nieudanej kalibracji. Dodatkowo ocenia się odsetek prób jednocześnie użytecznych i mieszczących się w progu aplikacyjnym. Osobę bez ważnego bloku wykazuje się w mianowniku dostępności. Przedziały ufności uwzględniają grupowanie pomiarów w osobach. Powtarzalność opisuje rozkład różnic między dwiema wizytami oraz ich granice zgodności dla uprzednio określonego wskaźnika na osobę. Zamrożenie estymandu, tolerancji aplikacyjnej i zasad obsługi jakości nastąpi przed zebraniem próby głównej.

\section{Wyniki}
Nie zebrano danych uczestników. Dwa syntetyczne przypadki z identycznym przesunięciem $100$ pikseli, przy skali $0{,}1$ mm na piksel i odległości $400$ mm, dały około $1{,}432^\circ$ przy osi oraz $1{,}406^\circ$ poza osią. Obliczenia sprawdzono względem rozwiązania analitycznego w płaszczyźnie. Ten wynik dokumentuje transformację geometryczną.

\section{Dyskusja i ograniczenia}
Błąd względem celu obejmuje także odchylenie uczestnika od instrukcji fiksacji. Analiza ograniczona do udanych prób może zaniżać obraz trudności w warunkach o dużej utracie śledzenia. Dlatego wynik warunkowy interpretowany jest razem z dostępnością oszacowań. Porównania z urządzeniem referencyjnym wymagają synchronizacji i określenia niepewności obu metod. Gotowość rekrutacyjna zależy od kwalifikacji E00, zamrożenia procedury E01 i uzyskania właściwej decyzji etycznej.
\bibliographystyle{apalike}
\bibliography{references}
\end{document}
